\documentclass[11pt,a4paper]{article}

\usepackage[utf8]{inputenc}
\usepackage[T1]{fontenc}
\usepackage{lmodern}
\usepackage[margin=2cm]{geometry}
\usepackage{amsmath}
\usepackage{booktabs}
\usepackage{enumitem}
\usepackage{titlesec}
\usepackage{hyperref}

\setlist{itemsep=1pt, topsep=2pt}
\setlength{\parskip}{3pt}
\setlength{\parindent}{0pt}
\titlespacing*{\section}{0pt}{8pt}{3pt}
\newcommand{\code}[1]{\texttt{#1}}

\begin{document}

\begin{center}
{\Large \textbf{Predicting the Success of Steam Games -- Week 2 Report}}\\[3pt]
Yohan ALLOUCHE \quad Arthur DISERBEAU \quad Jason ZOH \quad Lucas GREZES\\
\small ESILV A4 -- Machine Learning -- November 2026
\end{center}

\textbf{Goal.} We predict if a Steam game will be a success one year after its release:
\textbf{top 20\%} for the number of reviews in its first year \textbf{and} \textbf{at least 80\%} positive reviews.
We predict at $T = 0$ (release) and at $T = 7$, $30$ and $90$ days.
Modules: data (Yohan), target and splits (Arthur), features and EDA (Jason), baseline and evaluation (Lucas).

% =====================================================================
\section{Clean and preprocess the dataset (Yohan, Jason, Lucas)}

We use the Kozyriev reviews (about 41 M reviews, 50,872 games) and the FronkonGames metadata (125,855 games).
\begin{itemize}
  \item \textbf{Loading.} The 2 GB review file is read in chunks, with only the useful columns, and each chunk is filtered (keeping all chunks crashed the memory).
        A broken header in FronkonGames shifted every column; after the fix, 77.5\% of the Kozyriev games are found.
  \item \textbf{Sampling.} 10,000 games released in 2015--2021 (seed 42), sampled \textbf{by game}: sampling reviews would change the target.
  \item \textbf{Reviews.} \code{days} = review date $-$ release date; we remove the 6.45\% with \code{days} $< 0$ (early access). 6,833,248 reviews remain.
        Left join on \code{app\_id} (\code{validate="one\_to\_one"}): 7,797 games found in FronkonGames, 2,203 missing.
  \item \textbf{Features.} Log price, free, platforms, age, languages, release month and weekday, description length,
        genres / categories / tags as 0/1 columns (tags kept if in $\geq 1\%$ of the \textbf{training} games, 217 columns).
        For $T > 0$: number of reviews and share of positive reviews before $T$.
        Imputation and scaling are done in the model pipeline, fitted on the training set only.
\end{itemize}

% =====================================================================
\section{Analyze features, missing values and target distribution (Jason, Arthur)}

\textbf{Missing values.} FronkonGames columns are empty for about 22\% of the games, at the same rate in the three splits (21.6\% to 23.2\%).
The share of positive reviews is empty when a game has no review before $T$ (25--28\% at $T = 30$); it is unknown, so not replaced by 0.

\textbf{Distributions.} Prices and review counts are very skewed (median of 12 reviews in the first year, max 71,110), so we use log transforms.
Extreme values (e.g.\ 103 languages) are real games and are kept.
\textbf{Spearman correlations} (training set): reviews at $T = 30$ is the strongest (0.53); at release, the number of languages (0.23).

\textbf{Leakage audit.} A feature is kept only if it is known at prediction time. We removed rating, positive ratio, final price,
discount, DLC count and release year; \code{helpful} votes and \code{hours} played (read at collection time: 13.6\% of the
first-month reviews have more hours than the time since release); and 9 Steam categories added after release
(\emph{Remote Play Together}, created in 2019, appears on 10.5\% of the 2015 games).

\textbf{Target.} The threshold is \textbf{66 reviews}; about \textbf{1 game out of 7} is a success (imbalanced).
The 2,200 games with no review are kept as failures (no survivorship bias). The success rate drops to 8--13\% in 2017--2018, then 15--18\% from 2019.

% =====================================================================
\section{Define training, validation and test sets (Arthur)}

We split \textbf{by release date}: the model must predict future games from past ones, and a random split would be too optimistic.
Games of the same day stay in the same split. The top 20\% threshold is computed on the \textbf{training set only}.
All choices are made on validation; the \textbf{test set is used once}, at the end.

\begin{center}
\small
\begin{tabular}{l l r r}
\toprule
Split & Release dates & Games & Success rate \\
\midrule
Train (70\%) & 2015-01-01 to 2020-05-20 & 7,007 & 13.2\% \\
Validation (15\%) & 2020-05-21 to 2021-03-29 & 1,499 & 15.7\% \\
Test (15\%) & 2021-03-30 to 2021-12-31 & 1,494 & 14.9\% \\
\bottomrule
\end{tabular}
\end{center}

% =====================================================================
\section{Select and implement a baseline model (Lucas)}

One dataset per $T$ and a scikit-learn \code{Pipeline} fitted on the training set: median imputation and \code{StandardScaler} for numerical features,
\code{OneHotEncoder} for month and weekday (categories, not quantities), 0/1 features unchanged.
We compare a \code{DummyClassifier} (always ``failure'') with a \textbf{logistic regression} (\code{class\_weight="balanced"}).
It is our baseline because it clearly beats the Dummy at every $T$, handles the imbalance, gives probabilities, is fast,
and is \textbf{interpretable}: its coefficients revealed the Steam categories leak (\emph{Remote Play on Tablet}: 66.7\% of successes in train).

% =====================================================================
\section{Evaluate initial performance (Lucas, Arthur)}

Main metric: \textbf{PR-AUC} on the predicted probabilities, adapted to rare classes (random model = success rate, 0.157).
Accuracy is useless here: always saying ``failure'' gives about 85\%. Results on the validation set:

\begin{center}
\small
\begin{tabular}{l c c c c c c}
\toprule
Model & $T$ & PR-AUC & ROC-AUC & F1 & Precision & Recall \\
\midrule
Dummy & all & 0.157 & 0.500 & 0.000 & 0.000 & 0.000 \\
Logistic regression & 0  & 0.407 & 0.765 & 0.405 & 0.303 & 0.613 \\
Logistic regression & 7  & 0.743 & 0.919 & 0.621 & 0.505 & 0.804 \\
Logistic regression & 30 & 0.810 & 0.943 & 0.669 & 0.547 & 0.860 \\
Logistic regression & 90 & 0.850 & 0.959 & 0.705 & 0.586 & 0.885 \\
\bottomrule
\end{tabular}
\end{center}

\begin{itemize}
  \item At $T = 0$, the model is \textbf{2.6 times better than random}, but with many false alarms (precision 30\%).
  \item One week of reviews makes the PR-AUC jump from 0.407 to 0.743; the gain is smaller after.
        Part of it is mechanical (early reviews are also counted in the target): not leakage, but to keep in mind.
  \item Removing the leaking categories lowered the $T = 0$ score (0.435 to 0.411; 0.407 after the same-day split fix): the lower score is the honest one.
\end{itemize}

\textbf{Next steps.} Tune $C$, try random forest, gradient boosting and an MLP on validation, add bootstrap confidence intervals,
then evaluate the final model once on the test set.

\textbf{Use of AI.} We used an AI assistant for code and checked every answer; six errors were found and fixed
(memory crash, three leaks, month encoding, same-day games split between train and validation), see \code{AI\_AUDIT.md} and \code{PROMPT\_LOG.md}.

\end{document}
